\subsection{Model-size ladder}
\label{sec:ladder-method}

We train all five YOLO11 sizes under one protocol to measure the effect of model
size directly. This avoids a comparison based on published scaling numbers from
different protocols. YOLO11n, s, m, l, and x are trained on the Laboro split
for 60 epochs at 640 pixels with seed 0 and nominal batch 64. Their physical
batches, the numbers of images processed together in one training step, are 16,
16, 8, 8, and 4 respectively. These are the largest batches that fit in 24\,GB
of unified memory. Every checkpoint is exported to FP16 Core ML with
non-maximum suppression (NMS) included in the exported model, using confidence
0.25 and IoU 0.45. A detector proposes many overlapping boxes for the same
fruit; NMS keeps the most confident box and removes any other box that overlaps
it by more than the IoU threshold \cite{neubeck2006nms}. Including it in the exported
model means that the timed model is exactly the model that would be shipped in
the application.

We then score each size in three ways. Detection quality is measured by
mAP50--95 on the Laboro validation split. Counting quality is the best counting
MAE over a sweep of the confidence threshold, reported for all stages together
and for fully ripe fruit alone. Computational cost is measured by Core ML
package size, model load time, and inference latency on both devices. We report
the median and the p90, an upper-tail latency that captures slow frames as well
as typical ones.
